\chapter*{Abstract}
\addcontentsline{toc}{chapter}{Abstract}

Multi-agent reinforcement learning (MARL) increasingly powers autonomous systems,
from traffic control to distributed energy grids, yet a fundamental question
remains open: when does an explicit \emph{coordination mechanism} actually help, and
when do independent learners suffice? This dissertation answers that question on a
controlled, inherently cooperative benchmark --- traffic signal control on a
$2\times2$ grid of four intersections in SUMO, integrated with Ray RLlib.

A Multi-Agent Proximal Policy Optimization (MAPPO) controller, with a centralized
critic and a neighbour-coupled reward, was implemented and trained to stable
convergence. Under a \emph{fair, matched-clearance} evaluation --- in which every
controller pays the same realistic safety clearance --- the learned controller
decisively outperforms both the fixed-time and max-pressure heuristics; a
methodological finding of this work is that the heuristics' apparent strength was an
artefact of an unfair zero-clearance evaluation. The project's MAPPO was further
confirmed to match the canonical published recipe of Yu~et~al.\ (2021), establishing
it as a faithful, competitive implementation rather than a weak strawman.

The central result is a surprise. Compared against Independent PPO (IPPO) --- a fully
decentralized comparator differing only in its decentralized critic and purely local
reward --- MAPPO proves \emph{statistically indistinguishable} on every
traffic-performance metric across five seeds; the only reliable difference is that
the cooperative controller performs slightly more phase switches, with no
performance benefit. On a $2\times2$ network the \emph{coordination value}
(MAPPO~$-$~IPPO) is therefore approximately zero. The interpretation is not that
coordination is worthless but that the network is small enough that incentives are
already locally aligned, so independent learners recover coordinated behaviour
implicitly. This locates precisely where explicit coordination should begin to pay
off --- at larger network scales --- and motivates the future work of scaling to a
$5\times5$ grid and of extending the study to genuine social dilemmas, both placed
beyond this submission by compute constraints. The contribution is thus a careful
negative result together with a sharpened hypothesis: the value of cooperation in
MARL is not a fixed property of an algorithm but a function of scale.
